\subsection{Tolerances}\label{sec:theory_general_tolerances}

Suppose you buy 100 resistors all labeled as \(100 \Omega\)
with a tolerance of \(\pm 5\%\). This means that the actual
resistance value can vary between \(95 \Omega\) and \(105 \Omega\)
even though it is labeled as \(100 \Omega\):
\[
  R = 100\,\Omega\,(1 \pm 0.05) \hspace{0.5cm} \Rightarrow \hspace{0.5cm}
  R_{\text{min}} = 100 \cdot 0.95 = 95\,\Omega, \hspace{0.5cm} R_{\text{max}} = 100 \cdot 1.05 = 105\,\Omega
\]

In real-world resistors, capacitors, and inductors (and
actually a lot of other components) cannot be manufactured to
their exact specified values. This is due to natural
imperfections of process, material, and environment.

Refer to Section~\ref{sec:theory_general_probability} on how to deal
with this as a designer.

